\documentclass[11pt]{article}
\usepackage[margin=1in]{geometry}
\usepackage{amsmath,amssymb}
\usepackage{booktabs}
\usepackage{hyperref}
\usepackage{xcolor}
\usepackage{listings}
\usepackage{enumitem}

\lstset{
  basicstyle=\ttfamily\small,
  breaklines=true,
  frame=single,
  numbers=left,
  numberstyle=\tiny\color{gray},
  showstringspaces=false
}

\title{Experimental Plan: Remaining Experiments for\\
``Depth-Efficient Quantum Topological Data Analysis\\for Financial
Crash Early Warning''}
\author{Internal working document --- QCE 2026 submission}
\date{\today}

\begin{document}
\maketitle

\section*{Purpose}

This document enumerates the experiments that must be completed before
the paper's [XX] placeholders can be filled in. Each experiment is
tagged with its target paper section, required backend, estimated
runtime, and the specific numbers it must produce. All compute is on
BlueQubit (statevector) plus a local Qiskit Aer install for
density-matrix noise simulation.

\section{Experiment Inventory}

\begin{table}[h]
\centering
\small
\begin{tabular}{llllr}
\toprule
ID & Name & Backend & Target section & Priority \\
\midrule
E1  & Chronological split classification       & CPU only           & \S~Classification Performance         & P0 \\
E1b & OOD crisis evaluation (COVID, rate cycle)& CPU only           & \S~Classification Performance (OOD)   & P0 \\
E2  & Real-data loop recovery                  & BlueQubit SV       & \S~PCE Convergence Validation         & P0 \\
E3  & LGZ-QPE noise at $n_k\in\{8,16\}$        & Qiskit Aer DM      & \S~Noise Robustness                   & P0 \\
E4  & $\kappa=2$ vs $\kappa=3$ comparison      & BlueQubit SV       & \S~Encoding-Order Comparison          & P1 \\
E5  & Extended barren plateau ($n=14,16$)      & BlueQubit SV       & \S~Barren Plateau Analysis            & P1 \\
E6  & Hardware resource estimate               & Local transpile    & \S~Near-Term Hardware Resource Est.   & P1 \\
E7  & Codebase polish + hyperparameters.json   & --                 & Reproducibility                       & P2 \\
\bottomrule
\end{tabular}
\caption{Experiment inventory. P0 = blocks submission; P1 = strongly
recommended; P2 = polish.}
\end{table}

\section{Experiment Details}

\subsection{E1: Chronological Split Classification}

\paragraph*{Target.} Section~\emph{Classification Performance},
Table~\emph{Classification Performance: Chronological Split},
Figure~\emph{ROC Classification Split}.

\paragraph*{Backend.} CPU-only (pandas, scikit-learn,
matplotlib). No quantum compute.

\paragraph*{Inputs.} Existing
\texttt{3\_Results/verification\_results.json} (Betti curves per
window), \texttt{data/sp500\_2003\_2010.csv} (prices).

\paragraph*{Procedure.}
\begin{enumerate}[leftmargin=*]
  \item Load per-window $\beta_1$ values and window start dates.
  \item Compute 90-day forward drawdown per window; label as 1 if
        drawdown $>$ 10\%, else 0.
  \item Compute 30-day rolling volatility proxy per window.
  \item Chronological split: \texttt{train = windows starting in
        2003--2006}; \texttt{test = windows starting in 2007--2010}.
  \item Sweep thresholds on training split only; pick $F_1$-optimal
        $\beta_1$ threshold and volatility threshold.
  \item Evaluate fixed thresholds on held-out test split.
  \item Record precision, recall, $F_1$, ROC AUC for both classifiers
        on both splits.
\end{enumerate}

\paragraph*{Outputs.}
\begin{itemize}[leftmargin=*]
  \item \texttt{3\_Results/classification\_chronological.json} with
        all eight (classifier $\times$ split $\times$ metric) values,
        plus \texttt{threshold\_beta\_1}, \texttt{threshold\_vol},
        split counts, positive counts.
  \item \texttt{Figures/fig\_roc\_classification\_split.png} with
        four curves (train/test $\times$ $\beta_1$/volatility).
\end{itemize}

\paragraph*{Placeholders filled.} Table I in paper (eight metric
cells), prose values [XX] in the Classification Performance section
(precision, recall, $F_1$, two AUCs, positive-window count).

\paragraph*{Runtime estimate.} Minutes (CPU only).

\subsection{E1b: Out-of-Distribution Crisis Evaluation}

\paragraph*{Target.} Table~\emph{Out-of-Distribution Evaluation}.

\paragraph*{Backend.} CPU-only.

\paragraph*{Inputs.} Fresh S\&P 500 downloads:
2019-06-01 to 2020-12-31 (COVID), 2022-01-01 to 2023-06-30 (rate
cycle). Use \texttt{yfinance}.

\paragraph*{Procedure.}
\begin{enumerate}[leftmargin=*]
  \item Download prices; compute log returns.
  \item Apply the identical Takens pipeline ($\tau=13$, $m=4$, $W=500$,
        step 8). Note: COVID window is short, so expect $\sim$30--50
        windows per episode.
  \item Compute Betti curves via \texttt{ripser} for each window (use
        the same $\varepsilon^*=0.32$ applied to the 2003--2010 data).
  \item Apply the 2003--2006-trained $\beta_1$ threshold from E1.
        \emph{Do not re-tune.}
  \item Compute precision/recall/$F_1$/AUC for each episode
        separately.
\end{enumerate}

\paragraph*{Outputs.}
\texttt{3\_Results/classification\_ood.json}; numbers go into
Table~II of the paper.

\paragraph*{Placeholders filled.} Eight metric cells in the OOD
table.

\paragraph*{Runtime estimate.} Under an hour (data download +
\texttt{ripser}).

\subsection{E2: Real-Data Loop Recovery}

\paragraph*{Target.} Prose paragraph in \S~PCE Convergence
Validation (second paragraph currently containing seven [XX]
placeholders).

\paragraph*{Backend.} BlueQubit statevector (\texttt{cpu} device is
sufficient for $n \leq 16$).

\paragraph*{Goal.} Demonstrate that PCE-VQE + deflation recovers
loops from a \emph{real} market Laplacian, not just toy examples.
This closes the ``toy only'' gap that reviewers will flag.

\paragraph*{Procedure.}
\begin{enumerate}[leftmargin=*]
  \item Identify the sliding window with maximal full $\beta_1$ from
        the 2007--2010 data (window~167, $\beta_1 = 22$).
  \item Run \texttt{ripser} on the full point cloud of that window
        to obtain the $H_1$ persistence diagram.
  \item Sort 1-cycles by persistence length; select top
        $L \in \{3, 5, 10\}$ most persistent.
  \item Build a subsampled edge set that contains every edge
        participating in those top-$L$ cycles, plus a controlled
        number of surrounding edges to reach $n_k$ in
        $\{64, 128, 256\}$.
  \item Assemble $\Delta_1^{\mathrm{sub}}$; classically compute
        $\beta_1^{\mathrm{sub}}$ via \texttt{scipy.linalg.eigh}.
  \item Run PCE-VQE (\texttt{PCEVQE(kappa=2, backend='bluequbit')})
        with deflation for up to $\beta_1^{\mathrm{sub}} + 1$ rounds.
  \item Record how many null vectors are successfully recovered
        before the loss threshold is exceeded.
\end{enumerate}

\paragraph*{Cycle-to-edge mapping.} \texttt{ripser} does not return
cycle representatives by default. Options:
\begin{enumerate}[label=(\alph*),leftmargin=*]
  \item Use \texttt{ripser.representation} extension if available;
        returns chain representatives.
  \item Use \texttt{gudhi.RipsComplex} with
        \texttt{persistence\_intervals\_in\_dimension(1)} and
        \texttt{persistence\_pairs()} which yield simplex indices.
  \item \textbf{Fallback (recommended for paper deadline):} instead
        of exact representative extraction, greedily select edges
        with length in the persistence interval $[b, d]$ of each
        top-$L$ cycle. This preserves most cycles empirically even
        if a few are broken.
\end{enumerate}

\paragraph*{Outputs.}
\texttt{3\_Results/real\_data\_loop\_recovery.json} containing
\texttt{window\_idx}, \texttt{beta\_1\_full}, \texttt{nk\_full},
\texttt{nk\_sub}, \texttt{n\_qubits}, \texttt{beta\_1\_sub},
\texttt{beta\_1\_recovered}, \texttt{loss\_history}.

\paragraph*{Placeholders filled.} All seven [XX] values in the
``real market structure'' paragraph:
$\beta_1^{\text{full}}$, $n_k^{\text{full}}$, $n_k$, $n$,
$\beta_1^{\text{sub}}$, recovered count, target count.

\paragraph*{Runtime estimate.} 4--8 hours including debug time. Each
PCE-VQE run at $n = 16$ takes $\sim$10--30 minutes for full COBYLA
convergence with 2--3 deflation rounds.

\paragraph*{Success criterion.} \emph{Any} nonzero recovery count is
publishable. Even 1/5 is better than the current validation at
$\beta_1 = 0$ only. If nothing recovers, report that honestly and
discuss in Limitations.

\subsection{E3: LGZ-QPE Noise Simulation at Matched Scale}

\paragraph*{Target.} \S~Noise Robustness (prose values), updated
Figure~\emph{Noise Robustness v2}.

\paragraph*{Backend.} Qiskit Aer \texttt{AerSimulator(method=
'density\_matrix')}. Runs on CPU in seconds at $n \leq 8$ system
qubits + 4 ancilla.

\paragraph*{Why not BlueQubit.} Statevector cannot represent mixed
states under depolarizing noise. You need density-matrix simulation
\emph{or} trajectory averaging. Qiskit Aer's density-matrix backend
is the simplest and fastest path.

\paragraph*{Test cases.}
\begin{itemize}[leftmargin=*]
  \item $n_k = 8$, $\beta_1 = 1$: 4-cycle padded with isolated
        diagonal entries.
  \item $n_k = 16$, $\beta_1 = 2$: two disjoint 4-cycles.
\end{itemize}

\paragraph*{Procedure.}
\begin{enumerate}[leftmargin=*]
  \item Assemble each test Laplacian; verify $\beta_1$ classically.
  \item Pad to dimension $2^q$ (q=3 or 4) with
        $\lambda_{\max}/2$ on diagonal.
  \item Compute $U = \exp(i \Delta t)$ via matrix exponential;
        Trotterize for controlled-$U^{2^k}$ gates.
  \item Build QPE circuit with $q$ ancilla qubits, Hadamard layer,
        controlled-$U^{2^k}$ cascade, inverse QFT, measurement.
  \item For each error rate $p \in \{0, 10^{-4}, 5\times10^{-4},
        10^{-3}, 5\times10^{-3}, 10^{-2}\}$:
        \begin{itemize}
          \item Build a \texttt{NoiseModel} with one-qubit
                depolarizing error $p$ on all 1q gates and two-qubit
                depolarizing error $p$ on all CX gates.
          \item Run with 4096 shots starting from a uniform
                superposition over the system register (approximates
                the maximally-mixed initial state).
          \item Estimate $\tilde\beta_1 = 2^q \cdot P(\text{ancilla
                = }0)$.
          \item Accuracy $= 1 - |\tilde\beta_1 - \beta_1^{\text{true}}|
                / \max(1, \beta_1^{\text{true}})$ clipped to $[0,1]$.
          \item Repeat 5 trials; average.
        \end{itemize}
\end{enumerate}

\paragraph*{Known pitfall.} Building
\texttt{controlled-exp(iDt)} from scratch is nontrivial.
Recommended: use \texttt{qiskit.circuit.library.PhaseEstimation}
with a custom unitary via \texttt{.unitary()}. This bypasses explicit
Trotterization and gives you a valid (if not hardware-realistic) QPE
circuit for the purpose of noise benchmarking.

\paragraph*{Outputs.}
\texttt{3\_Results/lgz\_qpe\_noise.json} with nested dict
\texttt{\{nk: \{p: \{beta\_est, beta\_true, accuracy\}\}\}}.

\paragraph*{Placeholders filled.} Five [XX] values in Noise
Robustness prose plus new data points in figure.

\paragraph*{Runtime estimate.} 1--2 days including debug time.
Actual simulation is fast; the controlled-unitary construction is
the time sink.

\subsection{E4: $\kappa=2$ vs.\ $\kappa=3$ Comparison}

\paragraph*{Target.} \S~Encoding-Order Comparison,
Table~\emph{Quadratic vs.\ Cubic PCE}.

\paragraph*{Backend.} BlueQubit statevector.

\paragraph*{Procedure.}
\begin{enumerate}[leftmargin=*]
  \item Construct benchmark Laplacian: two disjoint 32-edge cycles
        (n\_k = 64, $\beta_1 = 2$). Verify classically.
  \item Run PCE-VQE with $\kappa = 2$, $n = 8$ qubits, HEA layers
        $L = 16$, COBYLA, $\mu = 5.0$, $\delta = 0.01$, max\_iter =
        500. Execute 2 deflation rounds.
  \item Run PCE-VQE with $\kappa = 3$, $n = 5$ qubits, HEA layers
        $L = 10$, identical optimizer settings. Execute 2 deflation
        rounds.
  \item Record per run: qubit count, transpiled circuit depth,
        parameter count, mean COBYLA iterations to convergence
        per deflation round, whether each round successfully
        recovered a null vector.
  \item Repeat each run 5 times with different random seeds; report
        means.
\end{enumerate}

\paragraph*{Outputs.} \texttt{3\_Results/kappa\_comparison.json}
with one entry per $\kappa$.

\paragraph*{Placeholders filled.} Eight [XX] in
Table~\emph{Quadratic vs.\ Cubic} plus three prose [XX] values.

\paragraph*{Runtime estimate.} 2--4 hours including BlueQubit queue
time.

\subsection{E5: Extended Barren Plateau Sweep}

\paragraph*{Target.} \S~Barren Plateau Analysis, updated
Figure~\emph{Barren Plateau v2}.

\paragraph*{Backend.} BlueQubit statevector. $n=16$ is the heaviest
single run; still fast on \texttt{gpu} device.

\paragraph*{Procedure.}
\begin{enumerate}[leftmargin=*]
  \item Modify existing barren-plateau script: change sweep to
        $n \in \{4, 6, 8, 10, 12, 14, 16\}$.
  \item For each $n$, 500 random parameter initializations uniformly
        in $[-\pi, \pi]$.
  \item For each initialization and each deflation round
        $j \in \{0, 1, 2\}$, compute gradient of Rayleigh-quotient
        loss w.r.t.\ a fixed parameter index $\ell$ via parameter
        shift.
  \item Record variance over the 500 samples.
  \item Linear regression in log--log space: slope and $R^2$.
  \item Also compute slope from $n=12$ to $n=16$ only, to assess
        whether growth persists at larger $n$.
\end{enumerate}

\paragraph*{Outputs.}
\texttt{3\_Results/barren\_plateau\_extended.json} with
\texttt{\{j: \{n: variance\}\}}, full-range slope, and
high-$n$-only slope.

\paragraph*{Placeholders filled.} Six [XX] in Barren Plateau
prose (three full-range slopes + three high-$n$ slopes).

\paragraph*{Runtime estimate.} 4--12 hours depending on BlueQubit
queue. The $n=16$ sweep is the bottleneck; can parallelize seeds.

\paragraph*{Interpretation caveat.} If high-$n$ slopes turn
\emph{negative}, this is important to report honestly: it would
mean the finite-size non-vanishing we see breaks down, consistent
with asymptotic vanishing. Either outcome is publishable.

\subsection{E6: Hardware Resource Estimate}

\paragraph*{Target.} \S~Near-Term Hardware Resource Estimate.

\paragraph*{Backend.} Local Qiskit transpile only; no simulation run.

\paragraph*{Procedure.}
\begin{enumerate}[leftmargin=*]
  \item Build $n = 16$, $L = 32$ HEA via
        \texttt{TwoLocal(16, ['ry','rz'], 'cx', 'linear', reps=32)}.
  \item Transpile to basis \texttt{[cx, sx, rz]} at
        \texttt{optimization\_level=3} with a representative
        coupling map (IBM Heron-like, linear strip of 16 qubits).
  \item Count post-transpile: CX count, 1-qubit gate count, circuit
        depth.
  \item Compute per-shot fidelity:
        $f = (1 - \epsilon_{2Q})^{n_{CX}} \cdot
             (1 - \epsilon_{1Q})^{n_{1Q}}$
        with $\epsilon_{2Q} = 3\times 10^{-3}$,
        $\epsilon_{1Q} = 2\times 10^{-4}$.
  \item Shot budget: $N_{\text{shots}} = N_{\text{basis}} \cdot
        G \cdot N_{\text{iter}} \cdot N_{\text{rounds}}$ with
        $N_{\text{basis}} = 10{,}000$, $G = 3$, $N_{\text{iter}} =
        200$, $N_{\text{rounds}} = 2$.
\end{enumerate}

\paragraph*{Outputs.}
\texttt{3\_Results/hardware\_estimate.json}.

\paragraph*{Placeholders filled.} Six [XX] values in the Hardware
Resource Estimate section.

\paragraph*{Runtime estimate.} 10 minutes.

\subsection{E7: Codebase Polish}

\paragraph*{Target.} Reproducibility, Data Availability.

\paragraph*{Tasks.}
\begin{itemize}[leftmargin=*]
  \item Create \texttt{3\_Results/hyperparameters.json} logging
        optimizer settings, PCE parameters, Takens parameters, VR
        threshold, window settings, crash-label definition.
  \item Resolve the ``max 8.8'' ambiguity: the 8.8 is the
        \emph{window-averaged maximum} of nonzero-entries-per-row,
        not a single-row max. Rename the JSON field to
        \texttt{mean\_of\_max\_nnz\_per\_row} and update the paper
        text (already changed in this revision).
  \item Add a one-page \texttt{REPRODUCE.md} in the repo pointing
        each paper figure/table to its generating script and input
        files.
  \item Ensure fresh-clone reproducibility: run every script from a
        clean venv following \texttt{REPRODUCE.md} and confirm all
        outputs match committed JSONs / figures to within random
        seed.
\end{itemize}

\paragraph*{Runtime estimate.} A few hours.

\section{Backend and Tooling Setup}

\subsection{BlueQubit}

\begin{itemize}[leftmargin=*]
  \item Install: \texttt{pip install bluequbit}.
  \item Token: export \texttt{BLUEQUBIT\_API\_TOKEN} in shell.
  \item Device selection: \texttt{bq.run(circuit, device='cpu')} for
        $n \leq 20$; \texttt{device='gpu'} for $n > 20$ or when
        running many parallel seeds.
  \item Thin wrapper (already planned):
        \texttt{2\_Quantum/bluequbit\_backend.py} exposing
        \texttt{statevector(qc)} and \texttt{expval(qc, pauli\_str)}
        so the core PCE-VQE class does not hard-code the backend.
\end{itemize}

\subsection{Local Qiskit Aer (for E3 only)}

\begin{itemize}[leftmargin=*]
  \item Install: \texttt{pip install qiskit qiskit-aer}.
  \item No cloud token needed.
  \item Density-matrix backend:
        \texttt{AerSimulator(method='density\_matrix',
        noise\_model=nm)}.
  \item At $n \leq 8$ system qubits the density-matrix state has
        $\leq 2^{16}$ complex entries; runs in seconds per shot
        on a laptop.
\end{itemize}

\subsection{Why Not Pure Statevector for Noise?}

Statevector simulation models pure states $|\psi\rangle$ under
unitary evolution. Depolarizing noise maps pure states to mixed
states $\rho = \sum_k p_k |\psi_k\rangle\langle\psi_k|$, which a
single statevector cannot represent. Two acceptable workarounds if
local Qiskit is unavailable:
\begin{enumerate}[leftmargin=*]
  \item \textbf{Trajectory sampling:} on each shot, stochastically
        inject Pauli errors after each gate according to the
        depolarizing channel, then statevector-simulate the noisy
        circuit. Average over many trajectories. Implementable on
        BlueQubit but requires a custom pass; not recommended for
        the deadline.
  \item \textbf{BlueQubit density-matrix (if available):} if
        BlueQubit exposes a density-matrix backend at submission
        time, use it. Check BlueQubit docs; if not, use local
        Aer.
\end{enumerate}

\section{Execution Timeline}

\begin{table}[h]
\centering
\small
\begin{tabular}{lllll}
\toprule
Day & Experiment & Hours & Parallelizable with & Blocking? \\
\midrule
1    & E1, E1b, E6, E7    & 4--6  & all others   & No \\
1--2 & E4                 & 4     & E5, E3       & No \\
2--4 & E5 (submit to BQ)  & 12    & E2, E3       & No \\
3--5 & E3 (local Qiskit)  & 16    & E2, E5       & No \\
4--7 & E2                 & 16    & debug-heavy  & \textbf{Yes} (P0) \\
7    & Paper updates + figure regens & 4 & --    & \\
\bottomrule
\end{tabular}
\caption{Recommended timeline. Total calendar time: one week of
focused effort. E2 is the highest-risk item.}
\end{table}

\section{Risk Register}

\begin{itemize}[leftmargin=*]
  \item \textbf{E2 cycle extraction.} \texttt{ripser} does not cleanly
        expose cycle representatives. Fallback: length-based edge
        selection (preserves most cycles empirically). Have both
        implementations ready.
  \item \textbf{E3 LGZ circuit construction.} Building
        controlled-$\exp(i\Delta t)$ correctly is the main pitfall.
        Use \texttt{qiskit.circuit.library.PhaseEstimation} with a
        unitary matrix passed directly; accept that this is not
        hardware-decomposable but is valid as a logical QPE for
        noise benchmarking.
  \item \textbf{E5 unexpected slopes.} If high-$n$ slopes are
        strongly negative, the paper's finite-size claim weakens.
        Mitigation: report exactly what is measured, weaken prose
        to ``finite-size non-vanishing within the tested range,
        with an inflection near $n = 12$.''
  \item \textbf{OOD data quality (E1b).} 2020 COVID is an abrupt
        single event; few windows will cover the drawdown. Expect
        noisy AUCs. Mitigation: also report raw $\beta_1$ and
        volatility trajectories as a sanity-check figure; do not
        over-interpret single-episode AUC.
\end{itemize}

\section{Placeholder Dependency Map}

Every [XX] in the paper is filled by exactly one experiment. Losing
an experiment means its [XX] values must be either removed (prose
rewritten) or labeled as future work. Nothing else in the paper
should be silently left blank at submission.

\begin{table}[h]
\centering
\small
\begin{tabular}{ll}
\toprule
Paper location & Experiment \\
\midrule
Abstract (held-out ROC AUC)                         & E1 \\
Section I, Contribution item 4 (prose)              & E1, E1b, E3, E4 \\
Section VII.B, real-data paragraph (7 values)       & E2 \\
Section VII.C, $\kappa$ table + prose (11 values)   & E4 \\
Section VII.D, noise prose (5 values)               & E3 \\
Section VII.E, barren plateau prose (6 slopes)      & E5 \\
Section VII.G, classification tables (12 + 8 cells) & E1, E1b \\
Section VII.H, hardware estimate (6 values)         & E6 \\
\bottomrule
\end{tabular}
\end{table}

\end{document}